% =====================================================================
% THESIS INTRODUCTION DRAFT — 2026-06-04
% Sections 1.1–1.5.  Copy into Overleaf chapter 1.
% TODO stubs are marked \todo{...} — verify before submission.
% =====================================================================

\chapter{Introduction}

% ─────────────────────────────────────────────────────────────────────
\subsection{Clinical Background: Oral Biofilm, Dysbiosis, and Peri-implantitis}
\label{sec:intro_clinical}

Dental implants have become the standard of care for replacing missing teeth, with
tens of millions placed annually worldwide~\cite{Kolenbrander2010OralMultispecies}.
Despite high short-term success rates, peri-implantitis --- a progressive
biofilm-driven inflammatory condition causing peri-implant bone resorption --- affects
an estimated 10--20\,\% of implants over a ten-year horizon and remains the leading
cause of late implant failure~\cite{Berglundh2018Peri, Schwarz2022PeriImplantitis}.
Because peri-implantitis is refractory to conventional antibiotics and lacks reliable
early biomarkers, understanding the microbial dynamics that trigger it is a prerequisite
for rational therapeutic intervention~\cite{Lamont2018Periodontal}.

Peri-implantitis is fundamentally a disease of \emph{dysbiosis}: a shift in the
composition and metabolic activity of the subgingival polymicrobial community from a
commensal-dominated, health-associated state to a pathogen-enriched, inflammation-driving
state~\cite{Marsh2003Dental, Hajishengallis2012Dysbiosis}.
The oral microbiome comprises more than 700 taxa interacting through mutualism,
competition, and metabolic cross-feeding~\cite{Kolenbrander2010OralMultispecies,
Dewhirst2010HOMD}; it is not a passive coloniser but an actively self-organised
community whose emergent state --- commensal or dysbiotic --- is determined by the
relative strengths of these inter-species interactions.
This bistability is the biological motivation for modelling the community with
attractor-state ODE dynamics rather than single-species growth curves.

The transition from health to dysbiosis follows a reproducible spatial and temporal
succession~\cite{Kolenbrander2006FnBridge}.
Early colonisers such as \textit{Streptococcus oralis}~(\textit{So}) and
\textit{Actinomyces naeslundii}~(\textit{An}) attach to the salivary pellicle and
create an oxidative, commensal-permissive microenvironment.
Their fermentative metabolism lowers local pH and reduces oxygen tension, conditioning
the substratum for secondary colonisers and eventually for strictly anaerobic
late-colonising pathogens~\cite{Lamont2018Periodontal}.

\textit{Fusobacterium nucleatum}~(\textit{Fn}) occupies the pivotal bridging niche
between early commensals and late pathogens~\cite{Signat2011Fn, Kolenbrander2006FnBridge}.
Through co-aggregation with both aerotolerant early colonisers and obligate anaerobes,
\textit{Fn} physically and metabolically connects the two flanks of the community:
it channels arginine-derived metabolites toward the anaerobic consortium and provides
adhesion scaffolding that recruits keystone pathogens to the developing
biofilm~\cite{Sakanaka2022FnArginine, Jiang2025FnReview}.
Accordingly, \textit{Fn} abundance is elevated in both periodontitis and
peri-implantitis lesions and is a consistent feature of the dysbiotic
attractor~\cite{Signat2011Fn}.

\textit{Porphyromonas gingivalis}~(\textit{Pg}) acts as a low-abundance keystone
pathogen that disproportionately reorganises the community and subverts host immunity
through gingipain protease activity and lipopolysaccharide
modifications~\cite{Hajishengallis2014Polymicrobial, Mysak2014Pg}.
Rather than causing disease by simple numerical dominance, \textit{Pg} modulates
complement signalling and suppresses neutrophil recruitment, creating a
nutrient-rich, inflammation-sustained niche in which the broader anaerobic consortium
can expand~\cite{Hajishengallis2012Dysbiosis, Lamont2018Periodontal}.
Recent full-length 16S and metatranscriptomic profiling of peri-implant submucosal
communities confirms that \textit{Pg}-associated guilds are enriched with increasing
probing depth and disease severity~\cite{joshi2025}, underscoring \textit{Pg}'s
keystone role in the clinical disease state.

To study these dynamics in a mechanistically controlled setting, Heine et al.~\cite{Heine2025PeriImplant}
developed a five-species in-vitro model comprising \textit{So}, \textit{An},
\textit{Veillonella} spp.~(\textit{Vei}), \textit{Fn}, and \textit{Pg}, grown in
the HOBIC (Hannover Oral BIofilm Chamber) flow reactor~\cite{Kommerein2018HOBIC}.
The five species were chosen to represent the four ecological roles
--- commensal early colonisers (\textit{So}, \textit{An}), metabolic bridge (\textit{Fn}),
anaerobic cross-feeder (\textit{Vei}), and keystone pathogen (\textit{Pg}) ---
while remaining small enough for rigorous Bayesian parameter
identification of the symmetric interaction matrix~$\mathbf{A}$.
Two community states (commensal~CS/CH, dysbiotic~DS/DH) $\times$ two cultivation
modes (static vs.\ HOBIC flow) define the four experimental conditions that anchor
Chapter~\ref{ch:heine}.

% ─────────────────────────────────────────────────────────────────────
\subsection{Computational Ecology --- gLV ODE Inference for Microbial Communities}
\label{sec:intro_ecol}

Quantitative models of microbial community dynamics cast inter-species interactions
as ordinary differential equations (ODEs).
The generalized Lotka--Volterra (gLV) model~\cite{Lotka1925Elements, Stein2013gLV}
describes the temporal evolution of $n$ guild relative abundances $\phi_i(t)$ as
\begin{equation}
  \dot{\phi}_i = \phi_i\!\left(b_i + \sum_{j=1}^{n} A_{ij}\,\phi_j\right),
  \quad i = 1,\ldots,n,
\end{equation}
where $b_i$ are intrinsic growth rates and $A_{ij}$ encodes the signed ecological
effect of guild $j$ on guild $i$.
A replicator formulation constrains the composition to the simplex
$\sum_i\phi_i = 1$~\cite{TaylorJonker1978Replicator, Hofbauer1998EGT}, which is
appropriate for 16S-derived relative-abundance data~\cite{Gloor2017Microbiome}.

\textbf{Inference challenge.}
Fitting $A \in \mathbb{R}^{n \times n}$ plus per-patient growth vectors from a cohort
of $p$ patients, each observed at $T = 3$ time points, yields $\mathcal{O}(n^2)$
free parameters against $p \cdot (T-1) \cdot n$ data constraints.
For the ten-patient Dieckow cohort with $n = 10$ guilds this ratio exceeds~5,
rendering the system severely underdetermined~\cite{Cao2017Identifiability, Castro2025scarce}.
Castro et al.~\cite{Castro2025scarce} show rigorously that under such conditions the
posterior over $A$ is structurally overparameterised: even optimal Bayesian inference
recovers at least one interaction sign incorrectly with near-certain probability for
$n \geq 4$.
Zapié{}n-Campos et al.~\cite{ZapienCampos2024Stochastic} further demonstrate that
matching higher distributional moments of abundance improves recovery of true
interaction parameters in the low-replicate regime directly analogous to our cohort.

\textbf{Hamilton replicator.}
To impose thermodynamic consistency, we adopt a variational (Hamilton-principle)
replicator~\cite{JunkerBalzani2021ExtendedHamilton, Klempt2024ContinuumGrowth} in which
the interaction matrix is required to be symmetric, $A = A^\top$.
Symmetry reduces the parameter count from $n^2$ to $n(n+1)/2$ and ensures that the
interaction network derives from a scalar potential, precluding limit cycles and
guaranteeing convergence to fixed-point attractors~\cite{Hofbauer1998EGT}.
The statistical-physics perspective of Pasqualini et al.~\cite{Pasqualini2024DisorGLV}
provides theoretical support for this bistable-attractor picture: disordered gLV
ecosystems exhibit a phase transition between stable (health-like) and unstable
(dysbiotic-like) regimes depending on the statistical structure of $A$.

\textbf{Bayesian inference and cross-validation.}
Model parameters are inferred via Transitional Markov Chain Monte Carlo
(TMCMC)~\cite{ChingChen2007TMCMC, Betz2016TMCMC}, a sequential Monte Carlo scheme
that anneals the likelihood from the prior to the full posterior and simultaneously
estimates the evidence for model comparison~\cite{Chopin2002SMC}.
A four-stage diagnostic workflow --- structural identifiability analysis, global
optimisation, local stability check, and predictive power assessment --- is applied
following Paredes-Vázquez et al.~\cite{Paredes2025ODECalib} to ensure the inferred
parameters are not artefacts of numerical overfit.
Generalisation is assessed by leave-one-patient-out cross-validation (LOO-CV): the
model is re-fitted on $p-1$ patients and trajectory predictions are evaluated on the
held-out patient, providing an out-of-sample test robust to the small cohort
size~\cite{Gelman2013BDA, Gibson2021MDSINE2}.

% ─────────────────────────────────────────────────────────────────────
\subsection{Metabolic Constraints as Ecological Priors --- AGORA and Sign Priors}
\label{sec:intro_prior}

Genome-scale metabolic models (GEMs) encode, for each organism, the full stoichiometry
of metabolic reactions and transport processes, enabling computation of feasible flux
distributions via flux balance analysis (FBA)~\cite{OrthThiele2010WhatIsFBA}.
Applied to multi-species communities, parsimonious FBA (pFBA)~\cite{Lewis2012pFBA}
predicts which metabolites are secreted by one guild and consumed by another, directly
encoding the metabolic basis of ecological interactions.
A cross-feeding relationship --- guild $j$ secretes compound $\alpha$ and guild $i$
takes it up --- implies a positive ecological effect of $j$ on $i$; conversely,
secretion of inhibitory compounds such as H\textsubscript{2}O\textsubscript{2} or
H\textsubscript{2}S implies antagonism.
These flux-derived cross-feeding signs thus translate into predicted \emph{signs}
of the off-diagonal entries of the gLV interaction matrix $A$.

The AGORA2 resource~\cite{Heinken2023AGORA2} provides curated GEMs for 7{,}302
human-associated microorganisms.
In the present work we selected guild-representative models for all ten Szafra\'nski
guilds, simulated each on an oral-fluid medium, and computed pairwise net cross-feeding
flows $F_{ij}$.
This yielded a third evidence layer (L3) for the sign prior, complementing the
experimental (L1) and predicted (L2) interaction evidence from Szafra\'nski
et al.~\cite{Szafranski2025CT}.
Together the three layers constrain 35 guild pairs, compared to 22 from L1+L2 alone.

Crucially, the prior penalises sign violations, not magnitude deviations.
The penalty
\begin{equation}
    \mathcal{P} = \sum_{(i,j):\, F_{ij} \neq 0}
    \frac{|F_{ij}|}{2\sigma^{2}}
    \!\left[\max\!\left(0,\,-\operatorname{sgn}(F_{ij})\,A_{ij}\right)\right]^{2},
    \quad \sigma = 0.15,
    \label{eq:sign_penalty}
\end{equation}
is zero whenever $A_{ij}$ and $F_{ij}$ share the same sign, and grows quadratically
only when they conflict.
This design is deliberate: FBA fluxes carry units of
mmol\,gDW\textsuperscript{-1}\,h\textsuperscript{-1} whereas gLV interaction strengths
are dimensionless compositional rates, and no calibrated unit conversion exists.
Attempts to use FBA magnitudes as quantitative priors (MacArthur consumer-resource
form~\cite{Marsland2019MicrobialCRM}) failed in validation, confirming that the
ecological information content of GEMs lies in the \emph{direction} of interaction,
not in its magnitude.
To our knowledge this is the first application of genome-scale metabolic models as
sign priors for ecological interaction inference in oral biofilm communities.

% ─────────────────────────────────────────────────────────────────────
\subsection{Spatial Biofilm Structure: FISH Microscopy and PDE Extension}
\label{sec:intro_spatial}

Oral biofilms on implant surfaces are not spatially homogeneous: early colonisers
such as \textit{Streptococcus oralis} and \textit{Actinomyces naeslundii} attach
to the titanium surface and form a structurally stable base layer, while
\textit{Fusobacterium nucleatum} acts as a physical and metabolic bridge species
that enables late colonisers such as \textit{Porphyromonas gingivalis} to
establish in the anaerobic, nutrient-depleted deep zone of the
biofilm~\cite{Kolenbrander2010OralMultispecies, Signat2011Fn}.
Fluorescence \textit{in situ} hybridisation (FISH) combined with confocal laser
scanning microscopy (CLSM) provides species-level, three-dimensional resolution
of this vertical organisation~\cite{Manders1993Colocalization}.
In the Heine HOBIC model system~\cite{Heine2025PeriImplant, Kommerein2018HOBIC},
FISH z-profiles for all five focal species were decoded from four fluorescence
channels --- including the dual-label signature of \textit{F.~nucleatum} ---
and averaged over axial depth bins to yield quantitative species abundance profiles
as a function of depth.

The zero-dimensional gLV model introduced in \S\ref{sec:intro_ecol} captures
temporal community dynamics but neglects spatial structure.
A natural extension replaces the purely temporal ODE with a one-dimensional
reaction--diffusion system in which each species abundance $\phi_i(z,t)$ obeys
\begin{equation}
  \partial_t \phi_i
    = \underbrace{f_i(\boldsymbol{\phi})}_{\text{gLV kinetics}}
    + \partial_z\!\left[\,D_i(\boldsymbol{\phi})\,\partial_z \phi_i\right],
\end{equation}
where $D_i(\boldsymbol{\phi})$ is a possibly cross-diffusion, density-dependent
transport coefficient encoding how the local community composition modulates
species-level motility or passive advection through the biofilm
matrix~\cite{Stewart1998Diffusion, Klempt2024ContinuumGrowth, PainterHillen2002VolumeFilling}.
Cross-diffusion terms couple the flux of species $i$ to the gradient of species
$j \neq i$, providing a PDE-level analogue of the ecological interaction structure
inferred from the ODE model~\cite{Freingruber2025CrossDiffusion}.
The experimentally measured FISH z-profiles serve as spatial calibration data:
by fitting the steady-state PDE solutions to the observed depth profiles, the
diffusion and cross-diffusion coefficients are estimated via an inverse problem,
extending the parameter inference framework from purely temporal to spatiotemporal
data.

% ─────────────────────────────────────────────────────────────────────
\subsection{Thesis Overview}
\label{sec:intro_overview}

This thesis is organised into five chapters following the three-pillar structure
introduced above.
Chapter~2 describes the datasets and shared preprocessing pipeline: the
Dieckow~2024 longitudinal 16S cohort used for primary inference, the
Botelho~2021 cohort used for cross-cohort sign validation, and the Heine~2025
\textit{in~vitro} time-series used for attractor identification and FISH imaging.
Chapter~3 (Pillar~I) develops the ecological inference framework: the generalized
Lotka--Volterra and Hamilton replicator ODEs, the three-layer AGORA2 metabolic
sign prior, and leave-one-patient-out cross-validation showing a 13\,\% reduction
in LOO-RMSE relative to unconstrained gLV.
Chapter~4 (Pillar~II) extends the 0D inference to a one- and two-dimensional
reaction--diffusion PDE, fits cross-diffusion coefficients to \textit{z}-resolved
CLSM FISH profiles, and demonstrates that cross-diffusion halves the spatial
residual versus standard Fickian diffusion.
Chapter~5 (Pillar~III) cross-validates the inferred interaction signs against
COMETS dynamic FBA simulations and FISH co-occurrence statistics.
Chapter~6 synthesises the findings, discusses limitations, and outlines future
directions.

\medskip
